% =====================================================================
%  Dokumentasi Fitur Aplikasi Lelang
%  Berisi penjelasan 5 sistem inti + referensi lokasi kode.
%  Salin isi file ini ke Overleaf. Jika ingin dipakai sebagai dokumen
%  berdiri sendiri, hapus tanda komentar pada blok preamble di bawah.
% =====================================================================

% ----------------------- PREAMBLE (opsional) -------------------------
% \documentclass[11pt,a4paper]{article}
% \usepackage[utf8]{inputenc}
% \usepackage[T1]{fontenc}
% \usepackage[bahasa]{babel}
% \usepackage{listings}
% \usepackage{xcolor}
% \usepackage{hyperref}
% \usepackage{geometry}
% \geometry{margin=2.5cm}
%
% \lstdefinestyle{code}{
%   basicstyle=\ttfamily\footnotesize,
%   breaklines=true,
%   numbers=left,
%   numberstyle=\tiny\color{gray},
%   keywordstyle=\color{blue},
%   commentstyle=\color{gray!70},
%   stringstyle=\color{teal},
%   frame=single,
%   rulecolor=\color{gray!40},
%   showstringspaces=false,
%   tabsize=2
% }
% \lstset{style=code}
%
% \begin{document}
% \tableofcontents
% \newpage
% ---------------------------------------------------------------------

\section{Menampilkan Card Post Lelang}
\label{sec:card-lelang}

Card lelang dirender oleh komponen \texttt{AuctionCard}
(\texttt{src/components/auction/AuctionCard.tsx}). Komponen ini menerima
sebuah objek \texttt{auction} beserta relasi opsional (penjual, kategori,
jumlah tawaran, dan rating penjual), lalu menampilkannya sebagai kartu yang
menautkan ke halaman detail \texttt{/auctions/[id]}.

Fitur penting pada card:
\begin{itemize}
  \item \textbf{Hitung mundur waktu (\emph{countdown})}. Sub-komponen
  \texttt{Countdown} memperbarui sisa waktu setiap 1 detik memakai
  \texttt{setInterval}. Jika sisa waktu $<$ 1 jam, badge berubah menjadi
  merah (status \emph{urgent}/``Hot'').
  \item \textbf{Pembaruan status real-time}. Card berlangganan (subscribe)
  perubahan baris tabel \texttt{auctions} melalui Supabase Realtime, sehingga
  status lelang (misalnya berubah menjadi ``Selesai'' atau ``Gagal'')
  langsung tercermin tanpa refresh.
  \item \textbf{Indikator kondisi barang} (mint, near mint, dsb.) dengan
  titik warna, harga tawaran tertinggi, harga ``Beli Langsung'' (jika ada),
  badge ``Gratis'' ongkir, jumlah tawaran, dan rating bintang penjual.
\end{itemize}

\begin{lstlisting}[language=JavaScript,caption={Hitung mundur \& deteksi status urgent (AuctionCard.tsx)}]
function Countdown({ endTime }: { endTime: string }) {
  const [time, setTime] = useState(null);
  useEffect(() => {
    setTime(getTimeRemaining(endTime));
    const t = setInterval(() => setTime(getTimeRemaining(endTime)), 1000);
    return () => clearInterval(t);
  }, [endTime]);

  if (!time) return <span>--:--</span>;
  if (time.expired) return <span>Berakhir</span>;
  const urgent = time.total < 3600; // < 1 jam -> tandai "Hot"
  return (
    <span>
      {time.days > 0
        ? `${time.days}h ${String(time.hours).padStart(2, "0")}j`
        : `${String(time.hours).padStart(2, "0")}:${String(time.minutes).padStart(2, "0")}:${String(time.seconds).padStart(2, "0")}`}
    </span>
  );
}
\end{lstlisting}

\begin{lstlisting}[language=JavaScript,caption={Subscribe perubahan status lelang secara real-time}]
useEffect(() => {
  if (status !== "active") return;
  const supabase = createClient();
  const channel = supabase
    .channel(`auction-card-${auction.id}`)
    .on("postgres_changes",
      { event: "UPDATE", schema: "public", table: "auctions",
        filter: `id=eq.${auction.id}` },
      (payload) => {
        const updated = payload.new;
        if (updated.status !== "active") {
          setStatus(updated.status);
          setWinnerId(updated.winner_id);
        }
      })
    .subscribe();
  return () => { supabase.removeChannel(channel); };
}, [auction.id, status]);
\end{lstlisting}

\paragraph{Referensi kode.}
\begin{itemize}
  \item \texttt{src/components/auction/AuctionCard.tsx} --- komponen card.
  \item \texttt{src/components/auction/AuctionBrowser.tsx} --- daftar/grid card.
  \item \texttt{src/app/api/auctions/route.ts} --- endpoint pengambilan data lelang.
\end{itemize}


\section{Sistem Kerja Transaksi Middleman}
\label{sec:middleman}

Middleman berperan sebagai pihak ketiga tepercaya (\emph{escrow}) yang menahan
dana pembeli sampai barang benar-benar diterima, baru kemudian melepaskannya ke
penjual. Alur ini melibatkan chat grup (GetStream) dan tabel \texttt{transactions}.

\subsection{Mengundang Middleman}
Di dalam percakapan, pengguna membuka \texttt{MiddlemanModal}
(\texttt{src/components/chat/MiddlemanModal.tsx}) yang mengambil daftar middleman
dari \texttt{/api/middlemen}, menyaring yang belum ada di percakapan, lalu
memanggil \texttt{onSelect(id, nama)}. Pemanggil kemudian menambahkan middleman
tersebut ke channel GetStream melalui \texttt{/api/stream/group}
(lihat Bagian~\ref{sec:getstream}), sehingga pembeli, penjual, dan middleman
berada dalam satu ruang percakapan yang sama.

\subsection{Membuat Transaksi (Escrow)}
Endpoint \texttt{POST /api/transactions}
(\texttt{src/app/api/transactions/route.ts}) memvalidasi bahwa: lelang sudah
selesai dan punya pemenang, pemohon adalah penjual/pembeli, belum ada transaksi
untuk lelang tersebut, dan (jika dipilih) middleman memang berstatus middleman.
Bila valid, sebuah baris transaksi dibuat dan lelang ditandai
\texttt{transaction\_confirmed = true}.

\begin{lstlisting}[language=JavaScript,caption={Pembuatan baris transaksi (transactions/route.ts)}]
const { data: transaction } = await supabase
  .from("transactions")
  .insert({
    auction_id,
    seller_id: auction.seller_id,
    buyer_id: auction.winner_id,
    middleman_id: middleman_id ?? null,
    amount: auction.current_price,
    tax_bearer: tax_bearer ?? "Buyer",
    stream_channel_id: stream_channel_id ?? null,
  })
  .select().single();

// Notifikasi ke middleman jika ada
if (middleman_id) {
  await createNotification(supabase, {
    user_id: middleman_id,
    type: "transaction_confirmed",
    auction_id,
  });
}
\end{lstlisting}

\subsection{Siklus Hidup Transaksi}
Status transaksi berjalan berurutan (lihat \texttt{TransactionCard.tsx}):

\begin{enumerate}
  \item \texttt{awaiting\_payment} --- Pembeli membayar; \textbf{middleman}
  menekan ``Konfirmasi Pembayaran Diterima''
  ($\rightarrow$ \texttt{/confirm-payment}).
  \item \texttt{awaiting\_shipment} --- \textbf{Penjual} mengunggah bukti
  pengiriman + nomor resi ($\rightarrow$ \texttt{/upload-shipping}).
  \item \texttt{shipped} --- \textbf{Pembeli} mengunggah/menekan konfirmasi
  penerimaan ($\rightarrow$ \texttt{/confirm-receipt}).
  \item \texttt{delivery\_confirmed} --- \textbf{Middleman} memverifikasi bukti
  lalu ``Lepaskan Dana ke Seller'' ($\rightarrow$ \texttt{/release-funds}).
  \item \texttt{completed} --- Transaksi selesai; statistik penjual \& pembeli
  bertambah, dan pembeli dapat memberi ulasan.
\end{enumerate}

Untuk \emph{direct deal} (tanpa middleman), langkah escrow dilewati dan
transaksi langsung ditampilkan sebagai selesai.

\begin{lstlisting}[language=JavaScript,caption={Pelepasan dana oleh middleman (release-funds/route.ts)}]
if (transaction.middleman_id !== user.id)
  return NextResponse.json({ error: "Hanya middleman yang bisa melepaskan dana" }, { status: 403 });

if (transaction.status !== "delivery_confirmed")
  return NextResponse.json({ error: `Status tidak valid` }, { status: 400 });

await supabase.from("transactions")
  .update({ status: "completed", updated_at: new Date().toISOString() })
  .eq("id", id);

// Tambah statistik & kirim notifikasi ke seller + buyer
await createNotification(supabase, { user_id: transaction.seller_id, type: "funds_released", auction_id: transaction.auction_id });
await createNotification(supabase, { user_id: transaction.buyer_id, type: "delivery_confirmed", auction_id: transaction.auction_id });
\end{lstlisting}

\paragraph{Referensi kode.}
\begin{itemize}
  \item \texttt{src/components/chat/MiddlemanModal.tsx} --- pemilihan middleman.
  \item \texttt{src/app/api/transactions/route.ts} --- membuat transaksi.
  \item \texttt{src/components/transaction/TransactionCard.tsx} --- UI \& stepper siklus.
  \item \texttt{src/app/api/transactions/[id]/confirm-payment} $|$ \texttt{upload-shipping} $|$ \texttt{confirm-receipt} $|$ \texttt{release-funds} --- transisi status.
  \item \texttt{supabase/migration-transactions.sql} --- skema tabel transaksi.
\end{itemize}


\section{Sistem Kerja Ulasan}
\label{sec:ulasan}

Ulasan hanya bisa diberikan oleh \textbf{pembeli} setelah transaksi berstatus
\texttt{completed}. Alur dimulai dari tombol ``Beri Ulasan Penjual'' di
\texttt{TransactionCard}, yang membuka \texttt{ReviewModal}
(\texttt{src/components/transaction/ReviewModal.tsx}).

\subsection{Penyimpanan Ulasan}
Endpoint \texttt{POST /api/transactions/[id]/review}
(\texttt{src/app/api/transactions/[id]/review/route.ts}) memvalidasi masukan
dengan Zod (rating bilangan bulat 1--5, komentar maksimum 1000 karakter), lalu
memastikan pemohon adalah pembeli dan transaksi sudah selesai. Ulasan disimpan
ke tabel \texttt{reviews} beserta \emph{snapshot} foto bukti penerimaan agar
dapat ditampilkan publik. Constraint \texttt{UNIQUE(transaction\_id)} mencegah
ulasan ganda (error kode \texttt{23505}).

\begin{lstlisting}[language=JavaScript,caption={Validasi \& penyimpanan ulasan (review/route.ts)}]
const reviewSchema = z.object({
  rating: z.number().int().min(1).max(5),
  comment: z.string().trim().max(1000).optional(),
});

if (transaction.buyer_id !== user.id)
  return NextResponse.json({ error: "Hanya pembeli yang bisa memberi ulasan" }, { status: 403 });
if (transaction.status !== "completed")
  return NextResponse.json({ error: "Ulasan hanya setelah transaksi selesai." }, { status: 400 });

const { error: insertErr } = await supabase.from("reviews").insert({
  transaction_id: transaction.id,
  auction_id: transaction.auction_id,
  seller_id: transaction.seller_id,
  buyer_id: user.id,
  rating,
  comment: comment || null,
  receipt_photo_url: transaction.receipt_photo_url ?? null,
});

if (insertErr?.code === "23505")
  return NextResponse.json({ error: "Anda sudah memberi ulasan." }, { status: 409 });
\end{lstlisting}

\subsection{Menampilkan Ulasan}
Komponen \texttt{SellerReviews}
(\texttt{src/components/profile/SellerReviews.tsx}) menampilkan rata-rata
rating, distribusi bintang 1--5 (dalam bentuk bar persentase), dan daftar
ulasan. Setiap ulasan menampilkan pembeli, waktu relatif, komentar, tautan ke
postingan pembelian, serta \emph{lightbox} foto bukti penerimaan.

\paragraph{Referensi kode.}
\begin{itemize}
  \item \texttt{src/components/transaction/ReviewModal.tsx} --- form ulasan.
  \item \texttt{src/app/api/transactions/[id]/review/route.ts} --- endpoint simpan ulasan.
  \item \texttt{src/components/profile/SellerReviews.tsx} --- tampilan ulasan penjual.
  \item \texttt{src/components/ui/StarRating.tsx} --- komponen bintang.
  \item \texttt{supabase/migration-seller-reviews.sql} --- skema tabel \texttt{reviews}.
\end{itemize}


\section{Sistem Kerja API GetStream (Chat)}
\label{sec:getstream}

Fitur chat memakai layanan \textbf{GetStream}. Prinsip keamanannya:
\texttt{API\_SECRET} \emph{hanya} boleh ada di server; browser cukup memakai
\texttt{API\_KEY} publik + \emph{token} pengguna.

\subsection{Server Client}
\texttt{src/lib/stream/client.ts} membuat koneksi GetStream sisi server dengan
hak admin (key + secret) sebagai \emph{singleton} agar tidak dibuat ulang tiap
request.

\begin{lstlisting}[language=JavaScript,caption={Server client GetStream (client.ts)}]
import { StreamChat } from "stream-chat";
let client: StreamChat | null = null;

export function getStreamServerClient(): StreamChat {
  if (!client) {
    client = StreamChat.getInstance(
      process.env.NEXT_PUBLIC_STREAM_API_KEY!, // key publik
      process.env.STREAM_API_SECRET!,          // secret (HANYA server)
    );
  }
  return client;
}
\end{lstlisting}

\subsection{Pengambilan Token}
Endpoint \texttt{POST /api/stream-token}
(\texttt{src/app/api/stream-token/route.ts}) adalah pintu masuk chat:
(1) memverifikasi sesi login Supabase, (2) mengambil profil,
(3) menyinkronkan profil ke GetStream (\texttt{upsertUser}), lalu
(4) membuat token JWT untuk pengguna. Token ditandatangani secara lokal memakai
secret sehingga cepat dan aman.

\begin{lstlisting}[language=JavaScript,caption={Membuat token GetStream (stream-token/route.ts)}]
const streamClient = getStreamServerClient();

// Sinkronkan profil ke Stream (fire-and-forget agar tidak memblok)
streamClient.upsertUser({
  id: profile.id,
  name: profile.name,
  username: profile.username,
  image: profile.avatar_url ?? undefined,
}).catch(console.error);

// createToken() = penandatanganan JWT lokal memakai API secret
const token = streamClient.createToken(profile.id);

return NextResponse.json({
  token,
  apiKey: process.env.NEXT_PUBLIC_STREAM_API_KEY!,
  userId: profile.id,
});
\end{lstlisting}

\subsection{Pembuatan Channel / Grup}
Endpoint \texttt{POST /api/stream/group}
(\texttt{src/app/api/stream/group/route.ts}) membuat channel dan menambahkan
anggota. Untuk fitur middleman, pembeli, penjual, dan middleman dimasukkan ke
satu channel. Keamanan: pemohon wajib termasuk dalam daftar member.

\begin{lstlisting}[language=JavaScript,caption={Membuat channel & menambah member (stream/group/route.ts)}]
if (!memberIds.includes(user.id))
  return NextResponse.json({ error: "Forbidden" }, { status: 403 });

const streamClient = getStreamServerClient();
const channel = streamClient.channel("messaging", channelId, {
  ...(name ? { name } : {}),
  created_by_id: user.id,
});

await channel.create();          // inisialisasi channel
await channel.addMembers(memberIds); // daftarkan anggota
\end{lstlisting}

\paragraph{Referensi kode.}
\begin{itemize}
  \item \texttt{src/lib/stream/client.ts} --- server client (key + secret).
  \item \texttt{src/app/api/stream-token/route.ts} --- pembuatan token.
  \item \texttt{src/app/api/stream/group/route.ts} --- pembuatan channel/grup.
  \item \texttt{src/components/chat/ChatWidget.tsx} --- koneksi browser (\texttt{connectUser}).
  \item \texttt{src/hooks/useChatUnread.ts} --- hitungan pesan belum dibaca.
\end{itemize}


\section{Sistem Kerja Notifikasi}
\label{sec:notifikasi}

Notifikasi disimpan di tabel \texttt{notifications} dan dikirim secara
real-time ke pengguna melalui Supabase Realtime.

\subsection{Membuat Notifikasi}
Helper \texttt{createNotification} (\texttt{src/lib/notifications.ts})
menyisipkan satu baris notifikasi. Tipe notifikasi bersifat terbatas (union
type), misalnya \texttt{auction\_won}, \texttt{transaction\_confirmed},
\texttt{funds\_released}, dan \texttt{auction\_approved}. Constraint
\texttt{UNIQUE(user\_id, type, auction\_id)} mencegah notifikasi duplikat.

\begin{lstlisting}[language=JavaScript,caption={Helper pembuat notifikasi (notifications.ts)}]
export async function createNotification(supabase, payload: {
  user_id: string;
  type: NotificationType;
  auction_id?: string;
}): Promise<void> {
  await supabase.from("notifications").insert({
    user_id: payload.user_id,
    type: payload.type,
    auction_id: payload.auction_id ?? null,
  });
  // Duplikat dicegah oleh UNIQUE (user_id, type, auction_id)
}
\end{lstlisting}

Helper ini dipanggil dari berbagai endpoint, misalnya saat lelang ditutup
(\texttt{cron/close-expired}), transaksi dikonfirmasi, dan dana dilepaskan
(Bagian~\ref{sec:middleman}).

\subsection{Mengambil Notifikasi}
Endpoint \texttt{GET /api/notifications}
(\texttt{src/app/api/notifications/route.ts}) mengambil 20 notifikasi terbaru
milik pengguna beserta relasi lelang, dan menghitung jumlah yang belum dibaca.

\begin{lstlisting}[language=JavaScript,caption={Pengambilan notifikasi (notifications/route.ts)}]
const { data: notifications } = await supabase
  .from("notifications")
  .select("*, auction:auctions(id, title)")
  .eq("user_id", user.id)
  .order("created_at", { ascending: false })
  .limit(20);

const unread_count = notifList.filter((n) => !n.is_read).length;
return NextResponse.json({ notifications: notifList, unread_count });
\end{lstlisting}

\subsection{Tampilan \& Real-time (Bell)}
Komponen \texttt{NotificationBell}
(\texttt{src/components/notifications/NotificationBell.tsx}) mengambil
notifikasi awal via REST, lalu \emph{subscribe} event \texttt{INSERT} pada tabel
\texttt{notifications} yang difilter per \texttt{user\_id}. Notifikasi baru
langsung muncul di dropdown dan badge angka bertambah. Saat dropdown dibuka,
notifikasi ditandai sudah dibaca melalui
\texttt{PATCH /api/notifications/read-all}.

\begin{lstlisting}[language=JavaScript,caption={Subscribe notifikasi real-time (NotificationBell.tsx)}]
const channel = supabase
  .channel(`notifications-bell-${channelId}-${userId}`)
  .on("postgres_changes",
    { event: "INSERT", schema: "public", table: "notifications",
      filter: `user_id=eq.${userId}` },
    (payload) => {
      const newNotif = payload.new;
      setNotifications((prev) => [newNotif, ...prev]);
      setUnreadCount((c) => c + 1);
    })
  .subscribe();
\end{lstlisting}

\paragraph{Referensi kode.}
\begin{itemize}
  \item \texttt{src/lib/notifications.ts} --- helper pembuat notifikasi.
  \item \texttt{src/app/api/notifications/route.ts} --- pengambilan notifikasi.
  \item \texttt{src/app/api/notifications/read-all/route.ts} --- tandai sudah dibaca.
  \item \texttt{src/components/notifications/NotificationBell.tsx} --- lonceng \& real-time.
  \item \texttt{src/components/notifications/NotificationItem.tsx} --- item notifikasi.
  \item \texttt{supabase/migration-notifications-ratings.sql} --- skema tabel notifikasi.
\end{itemize}

% \end{document}   % <- aktifkan jika memakai blok preamble di atas
